\documentclass[10pt]{letter}
\usepackage[a4paper, margin=0.85in]{geometry}
\usepackage[T1]{fontenc}
\setlength{\parskip}{5pt}
\signature{Harsh Bordekar \\ harshbordekar@gmail.com}
\begin{document}
\begin{letter}{Airbus Operations GmbH \\ Hamburg}
\opening{Dear Hiring Team,}

I am writing to apply for the Fatigue \& Damage Tolerance Engineer position in Hamburg. I am a structural analysis and simulation engineer with over seven years in aerospace structural integrity, spanning fatigue and damage tolerance of composite, hybrid, and metallic structures, and I would bring that depth to a role framed as an entry point into F\&DT.

Analyzing and validating structures against defined F\&DT methods is the core of my day-to-day work. During my PhD at DLR and TU Braunschweig I derived a certification-aligned failure criterion, a leakage envelope, that maps combined strain and load states to failure onset and defines an allowable-damage-limit philosophy consistent with regulatory boundary conditions. As a Research Intern in Metallic Fatigue and Damage Tolerance at DLR, I coupled defect measurements with numerical simulation to establish a defect-tolerance qualification methodology, translating fracture-mechanics and fatigue criteria into allowable-defect acceptance limits, the stress-analysis basis for accepting or rejecting a defective structure. I also validate theoretical damage models against measured data through Monte Carlo uncertainty quantification, and I work with progressive damage analysis of composites (CDM, PDA, CZM) alongside metallic fatigue assessment.

Beyond the analysis itself, I regularly work across design, materials, and stress interfaces: I provide sign-off-style technical review of externally and student-produced structural analysis work, define technical scope for that work, and communicate findings to design and project teams. I work daily with the FE solvers (Nastran/Patran, ABAQUS) and analytical stress methods (Niu, Bruhn) that sit alongside F\&DT sizing, build Python-based automation for analysis and reporting, and am genuinely interested in applying AI to improve engineering workflows.

To be transparent about fit: my acceptance-limit work has been for manufacturing-induced defects rather than in-service damage, so I have not designed or validated physical repair schemes, and I have not used Hyperworks; my FE tool experience is Nastran/Patran and ABAQUS. My German is currently at B1 level. I would rather name these directly than have them surface later.

I am motivated by the chance to bring rigorous, certification-oriented F\&DT thinking to Airbus, and I would welcome the opportunity to discuss how my background in damage-tolerance philosophy, model verification, and technical review can contribute to your team.

Thank you for your consideration. I look forward to discussing how I can contribute.

\closing{Sincerely,}
\end{letter}
\end{document}
